\section{Demo en tiempo real con un robot simulado}
\label{sec:demo}

Esta sección describe cómo se conecta el clasificador a un robot simulado en ROS\,2
\citep{macenski2022ros2} con \emph{turtlesim}. Usamos \emph{turtlesim} porque recibe las órdenes
de velocidad con la misma interfaz que los robots móviles reales (mensajes \code{Twist} en un
\emph{topic} \code{cmd\_vel}), así que el paso a un simulador como Gazebo, o a un robot físico, no
cambia el nodo de control.

\subsection{Diseño}

El sistema procesa cada fotograma de la cámara en cinco etapas (\cref{fig:pipeline}):
\begin{enumerate}
  \item \textbf{Captura} del fotograma con OpenCV.
  \item \textbf{Detección de la mano} con MediaPipe Hands \citep{zhang2020mediapipe}: la caja se
    obtiene como el rectángulo que envuelve los 21 puntos clave de la mano.
  \item \textbf{Recorte} con la misma función \code{crop\_hand} que se usó para generar los datos
    de entrenamiento (cuadrado, margen del 25\,\%, $128\times128$\,px) y normalización con las
    estadísticas de entrenamiento. Compartir esta función evita una fuente silenciosa de
    discrepancia entre entrenamiento e inferencia.
  \item \textbf{Clasificación} con GestureCNN, que devuelve una probabilidad por clase.
  \item \textbf{Filtrado y control}: una orden solo se emite si la probabilidad máxima supera
    0,7 \emph{y} es la clase mayoritaria en una ventana de los últimos 5--7 fotogramas. Si no se
    cumple, si no se detecta ninguna mano o si la clase es \cls{no\_gesture}, el nodo publica
    velocidad cero. Las velocidades se publican a 10\,Hz en \code{/turtle1/cmd\_vel}.
\end{enumerate}

El filtro de la última etapa responde al segundo requisito de la \cref{sec:introduccion}: un
fotograma aislado mal clasificado, por ejemplo durante la transición entre dos gestos, no debe
mover el robot. El umbral y la ventana intercambian seguridad por tiempo de reacción: a 30
fotogramas por segundo, una ventana de 7 fotogramas añade como mucho unos 0,2\,s de retardo.

\begin{figure}[t]
  \centering
  \begin{tikzpicture}[font=\footnotesize,
      st/.style={draw, rounded corners=2pt, minimum height=1.1cm, text width=1.62cm, align=center, fill=#1},
      arr/.style={-{Stealth[length=2mm]}, thick}]
    \node[st=gray!10]   (cam) {Cámara\\(OpenCV)};
    \node[st=green!12,  right=0.25cm of cam] (mp)  {MediaPipe\\Hands\\{\scriptsize caja de la mano}};
    \node[st=blue!10,   right=0.25cm of mp]  (cr)  {\code{crop\_hand}\\{\scriptsize $128^2$, margen 25\,\%}};
    \node[st=blue!20, text width=1.95cm, right=0.25cm of cr] (cnn) {GestureCNN\\{\scriptsize prob. por clase}};
    \node[st=orange!15, right=0.25cm of cnn] (fil) {Filtro\\{\scriptsize $p>0{,}7$ y voto 5--7}};
    \node[st=red!10,    right=0.25cm of fil] (ros) {Nodo ROS\,2\\{\scriptsize \code{Twist} a 10\,Hz}};
    \foreach \a/\b in {cam/mp,mp/cr,cr/cnn,cnn/fil,fil/ros} \draw[arr] (\a) -- (\b);
  \end{tikzpicture}
  \caption{Etapas de la demo en tiempo real. Las etapas 3 y 4 reutilizan exactamente el
    preprocesado y el modelo evaluados en este trabajo.}
  \label{fig:pipeline}
\end{figure}

\begin{table}[t]
  \centering
  \caption{Correspondencia entre gestos y órdenes para \emph{turtlesim}. Las órdenes de pinza
    no tienen equivalente directo en \emph{turtlesim}.}
  \label{tab:ordenes}
  \small
  \begin{tabular}{lll}
    \toprule
    Gesto & Orden & Acción en \emph{turtlesim} \\
    \midrule
    \cls{like}       & Avanzar               & \code{linear.x} $> 0$ \\
    \cls{dislike}    & Retroceder            & \code{linear.x} $< 0$ \\
    \cls{one}        & Girar a la izquierda  & \code{angular.z} $> 0$ \\
    \cls{peace}      & Girar a la derecha    & \code{angular.z} $< 0$ \\
    \cls{stop}       & Parar                 & velocidad cero \\
    \cls{fist}, \cls{palm} & Cerrar / abrir pinza & sin equivalente (ver texto) \\
    \cls{call}       & Volver al inicio      & servicio \code{teleport\_absolute} \\
    \cls{no\_gesture} & Ninguna              & velocidad cero \\
    \bottomrule
  \end{tabular}
\end{table}

\pendiente{indicar qué acción se asignó finalmente a \cls{fist} y \cls{palm} en \emph{turtlesim}
  (por ejemplo, subir y bajar el lápiz con el servicio \code{set\_pen}) y confirmar los valores
  finales del umbral, de la ventana de voto y de las velocidades lineal y angular.}

\subsection{Implementación y resultados}

\pendiente{describir la implementación del nodo \code{gesture\_controller} (versión de ROS\,2,
  entorno: Ubuntu nativo, máquina virtual o Docker) y cómo se integró el reconocedor con
  MediaPipe.}

\pendiente{tabla de latencias por etapa (captura, MediaPipe, recorte, CNN, filtro y extremo a
  extremo), con la media y el percentil 95 medidos con \code{time.perf\_counter} sobre al menos
  unos cientos de fotogramas, indicando el equipo usado. Como referencia, la etapa de
  clasificación por sí sola tarda \LatCpuFinal{} en CPU y \LatCudaFinal{} en GPU (mediana,
  \cref{tab:coste}).}

\pendiente{resultados cualitativos de la demo: si la tortuga obedece a cada gesto, activaciones
  falsas observadas (sobre todo durante las transiciones entre gestos), efecto del filtro y
  enlace al vídeo de la demostración.}
